\section{Experimental Setup}

\subsection{Datasets}

We evaluate SAQT on two vertical ancient-text recognition datasets, MTHv2 and HDRC. All experiments take cropped line-level or single-column images as input and predict the corresponding character sequence.

MTHv2 consists of three subsets and provides 105,579 vertical ancient-text single-column images in total~\cite{tang2020hrcenternet}. We unify the three subsets into a single recognition task for training and evaluation. During character-localization information learning, available character-level position annotations are additionally used as localization supervision after sorting boxes along the vertical direction.

HDRC comes from the ICDAR 2019 Historical Document Reading Challenge on Large Structured Chinese Family Records~\cite{saini2019hdrc}. The dataset provides PAGE-XML annotations. We construct line-level vertical recognition samples from text-line regions and Unicode transcriptions. To reduce the risk of samples from the same page appearing in both training and testing, HDRC is split by page groups. The train/validation/test splits contain 936/113/123 pages and 24,285/2,854/3,381 line samples. The charsets of MTHv2 and HDRC contain 6,727 and 4,017 characters, respectively.

\subsection{Baselines}

We compare SAQT with several scene text recognition baselines adapted to vertical input through MMOCR: CRNN, SVTR, ABINet, and SAR. Since these recognizers are designed for horizontal text by default, we rotate vertical single-column images at the input side so that their reading direction matches the expected model orientation. On MTHv2, the baselines include CRNN, SVTR-tiny, SVTR-small, SVTR-L, ABINet, and SAR. On HDRC, they include CRNN, SVTR-tiny, SVTR-L, ABINet, and SAR. All baseline outputs are converted to the same micro AR/CR evaluation protocol used for SAQT.

All baselines use dataset-specific MMOCR charsets and the same train/validation/test annotations. No external text corpus is used. Vertical images are rotated with \texttt{Rot90(k=3)}. CRNN uses height 32 and variable width up to 2048. SVTR-tiny, SVTR-small, and SVTR-L use the corresponding MMOCR configurations, with SVTR-L following the official 48$\times$160 geometry. ABINet uses 32$\times$128 input, and SAR uses padded images of height 48 and width 768. CRNN and SVTR variants are trained for 30 epochs; ABINet and SAR are trained for 20 epochs. Checkpoints are selected by validation \texttt{1-N.E.D}, and test predictions are then converted to AR/CR.

\subsection{Metrics}

We report accuracy rate (AR) and correct rate (CR). Let $S$, $D$, and $I$ denote substitution, deletion, and insertion errors, and let $N$ denote the number of ground-truth characters:
\begin{equation}
AR=1-\frac{S+D+I}{N},\qquad
CR=1-\frac{S+D}{N}.
\end{equation}

\subsection{Implementation Details}

SAQT uses a DINO-style detection transformer as the structure-query backbone. The visual backbone is ResNet-50. Both Transformer encoder and decoder contain 6 layers, with hidden dimension 256, 8 attention heads, and 900 queries. Character-localization information learning uses character-box supervision and DETR-style matching losses. The MTHv2 query-budget setting is trained for 6 epochs, with learning rate $1\times 10^{-4}$, backbone learning rate $1\times 10^{-5}$, batch size 2, and weight decay $1\times 10^{-4}$. The query-budget coefficient is 0.01.

For downstream recognition, query outputs are sorted by vertical position and organized as a CTC sequence. If the target charset differs from the reference charset associated with the localization-aware query representation, the classifier head is reconstructed before full-model recognition training. Shared characters inherit classifier weights, while new characters are initialized from unused classifier rows. Head reconstruction is trained for 3 epochs with learning rate $1\times 10^{-4}$. Full-model recognition training then continues for 12 epochs with learning rate $1\times 10^{-5}$. Both stages use AdamW, batch size 2, weight decay $1\times 10^{-4}$, and random seed 42. Model checkpoints are determined by development-set recognition results. Blank/nonblank calibration biases are determined on the development split and then fixed for testing; the test set is used only for final AR/CR reporting. The calibration introduces no language model or external text data. Experiments are conducted on NVIDIA RTX 3090 GPUs.
